\subsection{Why it is differentiable at forward cost}
\label{sec:theorem}

For the layer to be usable as a learned component it must be trainable. That it is,
at no asymptotic overhead, is a classical fact, stated once and built on.

\begin{remark}[Selected inversion is a gradient; its derivative has a sparse reverse schedule]
\label{rem:scaffold}
Write $P_{\pat}$ for orthogonal projection onto the on-pattern blocks. Under
the full Frobenius inner product on the symmetric matrix subspace, selected
inversion is the on-pattern gradient of the log-determinant,
$\selinv_{\pat}(A)=P_{\pat}\nabla_A\log\det A$
\citep{zhu2019sparse,dwyer1948symbolic}. Its derivative is therefore
self-adjoint after projection to that subspace. Reverse mode can be evaluated
by reversing the same elimination and Takahashi primitives, with no denser
pattern. By the cheap-gradient principle \citep{baur1983complexity}, its
arithmetic cost is a constant multiple of the forward's
$\Theta(Wb^3)$ work; the analytic implementation stores the factors and
selected blocks rather than an operation-by-operation autograd tape. For
non-symmetric $A$, write
$L_A:X\mapsto-\inv{A}X\inv{A}$ for the inversion derivative. Its adjoint is
$L_{A\T}$, not a selected inverse by itself. This transposes the
factor roles and requires independent lower and upper cotangents. The tree
specialization and projected-adjoint statement are in
Appendix~\ref{app:theorem}.
\end{remark}

Remark~\ref{rem:scaffold} is why the exact, global sparse solve of
\S\ref{sec:setup} can be used as a layer: embedded in a model and trained
end-to-end with the same asymptotic structural work as the forward pass, never forming the dense
inverse. From here the question is what the exactness of this layer does to
learning.
